\chapter{Abschlussprojekt}

\begin{lernziele}
Nach diesem Kapitel können Sie ein kleines, aber vollständiges MBSE-Projekt für einen ROS2-Roboter strukturieren.
\end{lernziele}

\section{Aufgabe}
Erstellen Sie ein vollständiges Modell des autonomen Indoor-Roboters. Das Modell soll Anforderungen, Struktur, Verhalten, Schnittstellen, ROS2-Mapping und Tests enthalten.

\section{Abzugebenes Modell}
Das Modell sollte folgende Inhalte enthalten:
\begin{itemize}
\item Kontextbeschreibung
\item Stakeholderliste
\item Mindestens zehn Anforderungen
\item Requirement Diagram
\item BDD mit Hauptsystemen
\item IBD für Perception oder Navigation
\item Aktivitätsdiagramm für Zielnavigation
\item Zustandsdiagramm für Betriebsmodi
\item Schnittstellentabelle
\item ROS2-Mapping-Tabelle
\item Mindestens fünf Testfälle
\item Traceability-Matrix
\end{itemize}

\section{Bewertungskriterien}
Ein gutes Modell ist verständlich, konsistent und nachvollziehbar. Es muss nicht maximal detailliert sein. Es muss zeigen, dass die Architektur aus Anforderungen abgeleitet wurde.

\section{Erweiterung}
Wer weitergehen möchte, kann eine einfache ROS2-Simulation ergänzen. Möglich sind ein simulierter Batterieknoten, ein Statusknoten und ein Dummy-Object-Detector.

\begin{uebungbox}
Erstellen Sie eine Traceability-Matrix mit den Spalten Anforderung, SysML-Block, ROS2-Node und Testfall.
\end{uebungbox}
